\section{Discussion}
\label{sec:discussion}

\subsection{Why HumanEval Training Dominates Both Benchmarks}

The embedding alignment framework explains the asymmetric specialization: HumanEval-only
training is more embedding-similar to both test benchmarks than MBPP-only training. But why?

HumanEval problems require writing functions that pass doctests --- a ``write correct Python code
that passes a specification'' skill that directly transfers to any test-case-based evaluation.
MBPP-only problems develop a narrower utility-script style (string manipulation, list operations)
that does not cultivate the same general function-writing competence. Both HumanEval+ and MBPP+
evaluate Python functions that must pass test cases; HumanEval's doctest-driven training format
more directly prepares the model for this evaluation structure.

An additional factor: the MBPP-only training condition used only 120 problems from the
sanitized split (vs the expected $\sim$374). This smaller set may have limited MBPP-only's
specialization potential. However, the direction (HumanEval-only $>$ MBPP-only on MBPP+) is
consistent across all 3 seeds, indicating a genuine effect.

\subsection{Equal-Mix Underperformance}

Equal-mix training achieves 10.2\% HumanEval+ pass@1 --- substantially below both single-source
conditions --- despite receiving the same token budget. Distributional dilution is the cause:
gradient updates are pulled in multiple distributional directions simultaneously, producing
weaker specialization than any focused single-source condition. This directly demonstrates that
source specificity dominates diversity at limited training scales.

\subsection{Scale and the $\eta^2$ Limitation}

The collapse of within-condition seed variance at 7B has a methodological implication: $\eta^2$ is
inappropriate for cross-scale effect-size comparison when seed sensitivity changes with scale.
Absolute condition spread (max mean $-$ min mean) is the correct primary metric. Future
scale-comparison studies in code SFT should report both.

\subsection{Limitations}

\textbf{MBPP+ data incomplete.} Full 4-condition MBPP+ evaluation was not saved in the H-E2 CSV
output (output path issue). MBPP+ claims are limited to the HumanEval-only vs MBPP-only
two-condition comparison from H-M1. Re-running evaluation on the existing 24 SFT checkpoints
would complete the full transfer matrix.

\textbf{Statistical fragility at n=4.} With four conditions, the minimum achievable Spearman
permutation p-value is 1/24 $\approx$ 0.042. Our $\rho$=1.0 achieves exactly this minimum. We interpret
this as ``highly consistent evidence for'' the alignment mechanism, not definitive proof.
Dual-encoder concordance (MiniLM $\rho$=0.8 in same direction) and the mechanistic plausibility
provide supporting evidence beyond the marginal p-value.

\textbf{Format normalization assumption.} MBPP-only SFT with native I/O example format (rather than
the standardized template) might produce stronger MBPP+ specialization, partially explaining the
absence of P2 inversion. Testing native format is an important future control.

\textbf{LeetCode training convergence.} LeetCode-only SFT achieved 0.0\% HumanEval+ pass@1 on 2/3 seeds
at 1.3B scale (seeds 42 and 123), with only seed 777 reaching 9.1\%. While the interpretation as
a source identity effect is plausible --- LeetCode's competitive-programming style is
distributionally divergent from HumanEval+ --- we did not separately verify training convergence
(loss curves) for LeetCode conditions. Future work should confirm LeetCode SFT training stability
at this scale to rule out training collapse as a confound.

\subsection{Broader Impact}

Benchmark comparisons should control for SFT source composition to enable valid attribution.
Reported pass@1 improvements may reflect source-benchmark alignment rather than architectural
advances. Positively, embedding cosine similarity between training candidates and target
benchmarks is a computationally cheap predictor of SFT performance rank, enabling principled
data curation before training.
